
\section{Preliminaries}

\subsection{Notation}
$\mathbb{L}^n$ ($\mathbb{L}_+^n$) denotes the space of $n \times n$ lower triangular (LT) matrices (with positive diagonal entries). $\mathbb{S}^n_+$ denotes the space of $n \times n$ symmetric positive definite  matrices. The Cholesky factor of a matrix \( P \in \mathbb{S}_+^{n} \) is defined as the unique matrix \( S \in \mathbb{L}_+^n \) such that \( P = S S^\top \) \cite[Theorem 4.2.7]{golubMatrixComputations2013}; we write \(S = P^{1/2}\). $\mathrm{diag}(M)$ denotes the vector of diagonal elements of matrix \( M \), and $\mathrm{diag}(a,b, \cdots)$ denotes the diagonal matrix with $a,b, \cdots$ on the diagonals. The linear matrix inequality $A \succ (\succeq) B$ denotes that $A - B $ is positive (semi)definite. \([N] := \{0, 1, \cdots, N\}\) for $N \in \N$.
We refer to the ``square root" of the covariance matrix \(P\) to refer to (not necessarily square) matrices \(L\) that satisfy \(P = L L^\top\). 
$\mathrm{tril}(\cdot)$ sets all but the LT part of a matrix to zero.
$\mathrm{vec}(\cdot)$ vectorizes a matrix.
$\mathrm{vectril}(\cdot)$ vectorizes the LT part of a matrix.

\subsection{Problem Formulation}
We address the finite-horizon discrete-time CS problem for a linear time-varying system with additive Gaussian noise. The goal is to find a state feedback controller that steers the state mean and covariance from their initial values to desired terminal values while minimizing a cost function and obeying CCs. Mathematically, the problem is formulated as follows.
\begin{subequations}
\label{eq:problem_statement}
\begin{align}
	& \mathrlap{\minimize_{\bm{v}, \bm{K}} \quad J = \sum_{k=0}^{N-1} J_k (x_k, u_k) }\\
	& \suchthat &&x_{k+1} = A_k x_k + B_k u_k + G_k w_k, \quad k \in [N{-}1] \\
	& && u_k = v_k + K_k (x_k - \E{x_k}), \quad k \in [N{-}1] \label{eq:feedback_control_law}\\
	& &&x_0 \sim \mathcal{N}(\mu_{\mathrm{init}}, P_{\mathrm{init}}) \\
	& &&\mu_{\mathrm{fin}} = \E{x_N}\\
	& &&P_{\mathrm{fin}} \succeq \E{(x_N - \E{x_N})(x_N - \E{x_N})^\top} \\
	& &&\mathbb{P}(x_k \in \mathcal{X}) \geq 1 - \epsilon_x, \quad k \in [N] \label{eq:state_chance_constraint}\\ 
	& &&\mathbb{P}(u_k \in \mathcal{U}) \geq 1 - \epsilon_u, \quad k \in [N{-}1] \label{eq:control_chance_constraint}
\end{align}
\end{subequations}
where \( x_k \in \mathbb{R}^n \) is the state vector, \( u_k \in \mathbb{R}^m \) is the control input, and \( w_k \in \mathbb{R}^n \) is standard Gaussian process noise. \( \bm{v} := \{v_k\}_{k=0}^{N-1}\) and \(\bm{K} := \{K_k\}_{k=0}^{N-1}\) are respectively the nominal control input and the state feedback gain to be designed, with \( v_k \in \mathbb{R}^m \) and \( K_k \in \mathbb{R}^{m \times n} \). Note that due to the affine feedback control law, the state and the control input are Gaussian at each time step.
The matrices \( A_k \), \( B_k \), and \( G_k \) are of appropriate dimensions.

We consider two types of cost functions:
\begin{equation} \label{eq:objective_function}
	J_k (x_k, u_k) = \begin{cases}
	\E{x_k^\top Q_k x_k + u_k^\top R_k u_k} \\
	\mathcal{Q}_{\norm{W_{k}^x x_k}_2}(p_J) +
	\mathcal{Q}_{\norm{W_{k}^u u_k}_2}(p_J)
	\end{cases}
\end{equation}
The former is a standard \textbf{expectation-of-quadratic} (EoQ) cost with \( Q_k \succeq 0 \) and \( R_k \succ 0 \).
The latter is a \textbf{quantile-of-norm} (QoN) cost, where the quantile function is defined as:
\begin{equation}
	\mathcal{Q}_{T}(\gamma) = \inf \{ t > 0 : \mathbb{P}(T \leq t) \geq \gamma \}.
\end{equation}
$W_k^x \in \mathbb{R}^{w_x \times n}$ and $W_k^u \in \mathbb{R}^{w_u \times m}$ are weighting matrices for the state and control, respectively. Such a cost formulation is useful for risk sensitive scenarios where the worst-case scenario (up to a specified probability) is of more interest than the expected value, such as in finance \cite{rockafellarOptimizationConditionalValueatrisk2000}, statistical machine learning \cite{eastwoodProbableDomainGeneralization2022}, and spacecraft mission design \cite{oguriChanceConstrainedControlSafe2024a}.

The sets \( \mathcal{X} \) and \( \mathcal{U} \) represent state and control feasible sets, respectively, with associated violation probabilities \( \epsilon_x \) and \( \epsilon_u \). Through some risk allocation procedure between each node, we assume that the CCs are given in one of two forms: a linear constraint or a norm constraint. For the linear constraint form, the CC is given by:
\begin{equation} \label{eq:linear_chance_constraint}
	\mathbb{P}(\alpha_j ^\top x_k \leq \beta_j) \geq 1 - p_j
\end{equation}
for each linear constraint defined by \( \alpha_j \in \mathbb{R}^n \) and \( \beta_j \in \mathbb{R} \), with violation probability \( p_j \).
For the norm constraint form, the CC is given by:
\begin{equation} \label{eq:norm_chance_constraint}
	\mathbb{P}(\| x_k \|_2 \leq \gamma) \geq 1 - p_\gamma
\end{equation}
for some scalar \( \gamma > 0 \) and violation probability \( p_\gamma \). The same forms apply to the control input \( u_k \).

For convenience, we make the following assumption which is reasonable in practical applications and also required in \cite{liuOptimalCovarianceSteering2025}:
\begin{assumption} \label{assump:positive_definite_state}
	The state covariance matrix remains strictly positive definite throughout the horizon.
\end{assumption}


\subsection{QR Decomposition}

The QR decomposition is a key element of propagating the square root of the covariance matrix, \textit{while keeping the square root to be in \(\mathbb{L}^n\)}. Various algorithms have been proposed for the time update of the square root Kalman filter, but many can be expressed as a QR decomposition \cite{tapleyStatisticalOrbitDetermination2004}.

\begin{lemma}[Existence and uniqueness of economy-size QR decomposition, Ch. 3.8 of \cite{ipsenNumericalMatrixAnalysis2009}]\label{lemma:existence_and_uniqueness_of_economy_size_qr_decomposition}
	Let \( M \in \mathbb{R}^{p \times q} \), \(p \geq q\).
	\begin{itemize}
		\item (Existence) There exist \(Q \in \mathbb{R}^{p \times q}\) and \(R \in \mathbb{R}^{q \times q}\) such that \(Q^\top Q = I_q\), \(R\) is upper triangular, and \(M = QR\).
		\item (Uniqueness) If \(M\) has full rank, i.e. \(\mathrm{rank}(M) = q\), then the matrices \(Q\) and \(R\) are unique up to a sign change of the diagonal entries of \(R\). If we require the diagonal entries of \( R \) to be positive, then the decomposition is unique.
	\end{itemize}
\end{lemma}
From here on, we will use the phrase ``QR decomposition'' to refer to the economy-size QR decomposition with positive diagonal entries.
\footnote{In MATLAB, this requires using $\mathrm{qr}(M, \text{``econ"})$. Then, to ensure uniqueness, if the function returns an \( R \) matrix with negative diagonal entries, multiply the corresponding columns of \( Q \) and rows of \( R \) by -1 to ensure positive diagonals.}

\begin{remark}
	\normalfont If \(M\) is full rank, then so is \(R\), hence \(R\) is invertible.
\end{remark}

\subsection{QR Decomposition Derivative}
In order to use this square root covariance propagation in an optimization framework, we utilize the derivative of the QR decomposition operation (i.e. the sensitivity of the $Q$ and $R$ matrices with respect to perturbations in the $M$ matrix). $\dd R$ is written as a function of \(Q\), \( R\), and  $\dd M$ \cite{deleeuwDifferentiatingQRDecomposition2023}:
\begin{equation}
	\dd R(\dd M, Q, R) = Q^\top (\dd M) - \mathrm{tril}(Q^\top (\dd M) R^{-1}) R.
\end{equation}
Using the Taylor expansion, we can approximate the QR decomposition as a linear function of perturbations in \( M \), around some reference \( \bar{M} = \bar{Q} \bar{R} \):
\begin{equation}
	\mathrm{qr}(M) \approx \bar{R} + \dd R(M - \bar{M}, \bar{Q}, \bar{R}).
\end{equation}
where \( \mathrm{qr}(\cdot) \) denotes the operation of extracting the upper triangular factor \( R \) from the QR decomposition.
\( \dd Q \) can also be computed as a function of \( \dd M\) \cite{deleeuwDifferentiatingQRDecomposition2023}; it is not used for this work, so we omit it here.